\section{Reproducible verification and the proof interfaces}
\label{app:reproduce}
The formalization has four components: physical block positions and their
multiplicities; the ordered quotient host; arithmetic classification and
repaired freeness; and transfer to matrices indexed by
$\operatorname{Fin}(n_h)$. The theorem
\code{end_to_end_main_theorem} combines these components with the
exact-count, arbitrary-edit distance and asymptotic results.

\subsection{Dependencies of the principal statements}
\begin{center}\small
\begin{tabular}{p{43mm}p{54mm}p{37mm}}
\toprule
Mathematical assertion&Lean declaration or module&Role\\\midrule
Arithmetic original classification&\code{Arithmetic.classification}&Copy exhaustion\\
Repaired freeness&\code{Arithmetic.repaired_no_copy}&No new copies\\
Encoding of physical copies&\code{Arithmetic.encode_copy}&Arithmetic soundness\\
Terminal decoding&\code{terminal_decodes}&Physical terminal forms\\
Numeric matrix order&\code{orderedPosition_order_iff}&Order transfer\\
Protected sampling&\code{XorSampling.lean}&Arbitrary edits\\
Finite avoidance&\code{FiniteAvoidance.lean}&Untouched witness\\
Copy multiplicity&\code{CopyFibres}, \code{TerminalCount}&Exact count\\
Edit distance&\code{EditDistance.lean}&Exact minimum\\
Asymptotics&\code{Limits.lean}&Removal failure\\
\bottomrule
\end{tabular}
\end{center}

The arithmetic classification and repair derivations contain 266 and
243 source segments, respectively. The encoding checks all eight
validity conditions, twelve axis-order comparisons and sixteen pattern
entries. The quotient-to-physical correspondence and the numeric-order
bijections complete the passage to the ordinary indexed matrix.
The proof segments and every imported Lean module are distributed as
source files in the accompanying verification archive.

\subsection{Rebuild and axiom audit}
The project uses Lean 4.24.0 and pins Mathlib to revision
\begin{center}\code{f897ebcf72cd16f89ab4577d0c826cd14afaafc7}.\end{center}
From the \code{formal_verification/} directory, execute
\begin{quote}\ttfamily\small
bash verify\_all.sh
\end{quote}
The script retrieves the pinned dependency cache, builds the primary
library, checks the main theorem and its axiom audit, then rebuilds the
isolated structural-extension and exact-transversal projects. The main
library is complete with all arithmetic certificates, which are needed
for the terminal classification and canonical freeness results.
The compiled dependency audits permit only \code{propext},
\code{Classical.choice} and \code{Quot.sound}.

The archived primary build record reports a compilation of
550 project modules, with 116,338 declarations (93,189 theorems)
in the strict declaration audit. The independent extension-core
build record covers 34 library modules and 1,570 declarations
(991 theorems). The exact-transversal supplement records 49 declarations,
including 32 theorems, and derives terminal exhaustion from the primary
global classification theorem. The corresponding run reports and source
identities are retained under \code{evidence/} and in the supplemental
project directories. The archived records identify the source and environment of earlier
compilations; the command above performs a new build on the recipient's machine.

\subsection{Mathematical cross-checks}
A direct check of the construction first verifies block-constant entries
and the four increasing forcing frames. The encoding implication then
maps every physical copy to the arithmetic host. The classification
certificate excludes all nonterminal copies; the repair certificate
excludes every copy after the canonical edit. Finally, weighted terminal
multiplicities give the exact number of original copies. The main
Lean theorem assembles these facts with the arbitrary-edit lower bound
and the asymptotic comparison.
